\section{All parabolic intersections in an affine period quotient}
\label{sec:affine-parabolics}

The finite adjacent subgroups in the polyhedral applications have faithful
linear images. An arbitrary parabolic need not have that property: it can
contain translations. We therefore state the intersection calculation before
making the faithful-image specialization. It uses the actual affine group,
whose identification with a presented group is the separate obligation in
AL.2.

Let $E=\langle r_i:i\in I\rangle\leq\mathbb Z^d\rtimes Q$, where $I$ and
the linear image $Q$ are finite. Multiplication is
$(q,u)(q',v)=(qq',u+qv)$. Write $L\leq\mathbb Z^d$ for the translation
kernel of $E$. For $J\subseteq I$ put
\[
 E_J=\langle r_i:i\in J\rangle,\qquad
 Q_J=\operatorname{im}(E_J\to Q),\qquad
 K_J=\{v:(1,v)\in E_J\}\leq L.
\]
Choose, for each $q\in Q_J$, an element $(q,u_J(q))\in E_J$, with
$u_J(1)=0$. These are finite choices of affine word lifts, not assumptions
about a splitting of $E_J$. The entire fiber over $q$ is
\begin{equation}
 \{(q,u):u\in u_J(q)+K_J\}.\label{eq:affine-fiber}
\end{equation}
Indeed the quotient of two elements with the same linear part is their
translation difference; conversely left multiplication by $(1,k)$ adds any
$k\in K_J$ to the translation part.

Let $M\leq L$ be $Q$-invariant. It is then a normal translation subgroup of
$E$. Write $\overline E=E/M$ and $\overline E_J$ for the image of $E_J$.
The subgroup $M$ can have any rank, including zero.

\begin{theorem}[AA.ParabolicPeriods]\label{thm:all-periods}
The marking of $E/M$ has the intersection property if and only if, for
every $J,T\subseteq I$, with $C=J\cap T$, both conditions hold:
\begin{align}
 (K_J+M)\cap(K_T+M)&=K_C+M,\label{eq:period-kernels}\\
 u_J(q)-u_T(q)&\notin K_J+K_T+M
 \quad(q\in(Q_J\cap Q_T)\setminus Q_C).\label{eq:period-exclusions}
\end{align}
More generally, a common linear part $q\in Q_J\cap Q_T$ occurs in
$\overline E_J\cap\overline E_T$ exactly when
\[
 u_J(q)-u_T(q)\in K_J+K_T+M.
\]
When it occurs, its translation fiber is a coset of
$((K_J+M)\cap(K_T+M))/M$.
\end{theorem}
\begin{proof}
Modulo $M$, equation~\eqref{eq:affine-fiber} becomes the coset
$(u_J(q)+K_J+M)/M$. Equality in $E/M$ preserves the linear part, since
$M$ consists of translations. Two cosets of additive subgroups intersect
if and only if the difference of their representatives belongs to the sum
of those subgroups. When nonempty, their intersection is a coset of the
intersection of the subgroups. This proves the two fiber assertions.

For $q\in Q_C$, a lift $(q,u_C(q))\in E_C$ belongs to both $E_J$ and
$E_T$. It can therefore be used as a common representative of their two
translation cosets. The intersection fiber is
\[
 \bigl(u_C(q)+(K_J+M)\cap(K_T+M)\bigr)/M,
\]
whereas the fiber of $\overline E_C$ is
$(u_C(q)+K_C+M)/M$. Condition~\eqref{eq:period-kernels} makes them equal.
Condition~\eqref{eq:period-exclusions} excludes every other possible linear
part. Together these prove
$\overline E_J\cap\overline E_T=\overline E_C$.

Conversely that equality, at the identity linear part, implies
\eqref{eq:period-kernels}: the relevant additive groups all contain $M$, so
equality of their quotients by $M$ is equality of the groups themselves.
For a linear part outside $Q_C$ its intersection fiber must be empty,
which gives~\eqref{eq:period-exclusions}. Applying this to every pair of
index sets proves the equivalence. Changing any chosen lift adds an element
of $K_J$ or $K_T$ and does not change the membership test.
\end{proof}

All data in this criterion can be obtained by finite exact calculations
for a fixed period lattice. Enumerate $Q_J$ from the linear generator
matrices, retain a word lift of each element, and apply AL.1 to $E_J$ and
$Q_J$. Its normalized Schreier translations generate exactly $K_J$.
Integer lattice sum, intersection and membership calculations then decide
\eqref{eq:period-kernels}--\eqref{eq:period-exclusions}. No finiteness of
$E_J$ and no injectivity of $E_J\to Q_J$ is assumed. For a parameterized
family $M_m$ the same equations remain exact, but an all-parameter solution
of them still needs a proof; finitely many successful periods do not
establish it.

The first condition cannot be omitted when parabolics contain translations.
Take $E=\mathbb Z^2$ with markers $e_1,e_2$ and trivial $Q$. For
$M=\langle e_1-e_2,2e_2\rangle$, the two distinct rank-one parabolics have
the same nontrivial image in $E/M\cong C_2$, so IP fails. Here
$K_{\{1\}}+M=K_{\{2\}}+M=\mathbb Z^2$, whereas $K_\varnothing+M=M$.
There is no nonidentity linear part to detect this failure. By contrast,
$M=2\mathbb Z^2$ gives the two distinct coordinate factors of $C_2^2$ and
satisfies the criterion. These examples concern general marked groups;
the theorem does not require involutory markers.

This statement extends the previously proved rank-three period argument to
all marked parabolics. Its proof is the classical intersection formula for
cosets of additive subgroups. The extension and the explicit hypothesis
separation are part of this edition; no claim of historical priority is
made. Existing formal period certificates cover the previously stated
faithful-parabolic monotonicity, not this full arbitrary-parabolic theorem.
